% Plantilla mínima para TFG (TFG / Trabajo Fin de Grado)
% Archivo: Info_TFG.tex
% Cómo usar: editar título/autor y dividir en capítulos con \include o \input si se desea.
% Compilar (recomendado): usar latexmk o pdflatex + biber/bibtex según bibliografía.
% Ejemplo PowerShell (TeX Live / MiKTeX instalado):
%   latexmk -pdf Info_TFG.tex
% o, paso a paso:
%   pdflatex Info_TFG.tex; biber Info_TFG; pdflatex Info_TFG.tex; pdflatex Info_TFG.tex

\documentclass[12pt,a4paper,oneside]{report}

% Codificación y idioma
\usepackage[T1]{fontenc}
\usepackage[utf8]{inputenc}
\usepackage[spanish]{babel}

% Márgenes y microtipografía
\usepackage[a4paper,margin=2.5cm]{geometry}
\usepackage{microtype}

% Gráficos, tablas y matemáticas
\usepackage{graphicx}
\usepackage{caption}
\usepackage{subcaption}
\usepackage{booktabs}
\usepackage{amsmath,amssymb}

% Hipervínculos
\usepackage[hidelinks]{hyperref}

% Numeración de secciones sin dependencia de capítulos (1,2,3... en vez de 0.1,0.2)
\usepackage{chngcntr}
\counterwithout{section}{chapter}

% Configuración del título/portada (personaliza según la normativa de tu universidad)
\title{Comparación de Nariz Electrónica con olfato de Agentes Caninos para la detección de Drogas}
\author{Raúl Martín-Romo Sánchez\\
	\textit{Grado en Ingeniería Informática del Software}\\
	\textbf{Tutor:} Pablo García Rodríguez\\
    \textbf{Co-tutor:} Jesús Lozano Rogado}
\date{\today}

\begin{document}

% Portada en la misma primera página (sin página separada)
\begingroup
\begin{center}
	{\LARGE \textbf{Comparación de Nariz Electrónica con olfato de Agentes Caninos para la detección de Drogas}\par}
	\vspace{1em}
	{\large Raúl Martín-Romo Sánchez\par}
	{\small Grado en Ingeniería Informática del Software\par}
	\vspace{0.5em}
	{\small \textbf{Tutor:} Pablo García Rodríguez \quad \textbf{Co-tutor:} Jesús Lozano Rogado\par}
	\vspace{1em}
	{\small \today\par}
\end{center}
\bigskip
\endgroup

% Capítulos de ejemplo
\section{Introduccion}
Hablar sobre la idea del proyecto y la motivación.

\section{Alcane y objetivos}
Describir los objetivos principales.

\section{Estado del arte}
Exponer la inforamación recogida acerca de las narices electrónicas, los sensores y algunos estudios de los diferentes artículos. Además hablar del cuerpo canino de la guardia civil, las funciones que desempeñan y como trabajan.

\section{Metodología y herramientas}
Planificación del proyecto con metodología Scrum, explicar los diferentes sprints y tareas a realizar añadiendo el diagrama de Gantt.
Explicar que materiales se han usado y para qué.

\section{Implementacion, desarrollo y pruebas comparativas}

\subsection{Entrenamiento de la e-nose}
Explicar como se ha creado el dataset y como se ha entradado la red neuronal.

\subsection{Desarrollo de aplicación Bluetooth}
Describir el desarrollo de la aplicación Bluetooth para la comunicación entre la nariz electrónica y el smartphone.

\subsection{Pruebas de comparación}
Describir que tipos de pruebas van a hacer.

\section{Resultados}
Descripción de los resultados esperados y de los datos obtenidos.
Discusión de los resultados obtenidos en comparación con los esperados.

\section{Conclusiones y futuras mejoras}
Describir la coclusiones sobre los resultados obtenidos, las posibles aplicaciones de la nariz electrónica y las futuras mejoras que se podrían implementar.

\end{document}

